\documentclass[12pt]{article}

\usepackage[
  a4paper,
  margin=1in,
  ]{geometry}

\usepackage[T1]{fontenc}

\PassOptionsToPackage{
  defaults=hu-min,
  labelenums=unchanged,
  postdescription=unchanged,
  mathfactorial=fix,
  }{magyar.ldf}

\usepackage[english,magyar]{babel}

\footnotestyle{rule=fourth}

\usepackage{
  xcolor,
  lmodern,
  upquote,
  enumitem,
  longtable,
  caption,
  tcolorbox,
  hyperref,
  }

\hypersetup{
  colorlinks,
  allcolors=blue,
  pdfstartview=FitH,
  bookmarksnumbered,
  }

\captionsetup{
  labelfont=bf,
  font=small,
%  width=10cm,
  labelsep=newline,
  justification=centering,
  }

\newlist{optionlist}{description}{3}
\setlist[optionlist]{nosep}
\setlist[description]{parsep=0pt}
\setlist[enumerate]{noitemsep}

\newcommand\param[1]{{\normalfont\slshape$\langle$#1$\rangle$}}

\renewcommand\descriptionlabel[1]{%
  \hspace{\labelsep}\normalfont\ttfamily\bfseries\hunotpvaluefont#1}

\def\hunotpvaluefont{}

\NewDocumentEnvironment{hunopt}{ m m o }{%
  \item[#1=\param{érték}]
  \hfill
  (alapérték: \textbf{\texttt{\color{gray}#2}})\\
  \IfValueT{#3}{#3 }%
  A lehetséges értékek:
  \begin{optionlist}
  \def\hunotpvaluefont{\color{gray}}
  }{
  \end{optionlist}
  }

\NewDocumentEnvironment{hunopt*}{ m m }{%
  \item[#1=\param{érték}]
  \hfill
  (alapérték: \textbf{\texttt{\color{gray}#2}})\\
  }{}

\NewDocumentEnvironment{footstyle}{ m }{%
  \item[#1=\param{érték}]
  A lehetséges értékek:
  \begin{optionlist}
  \def\hunotpvaluefont{\color{gray}}
  }{
  \end{optionlist}
  }

\tcbuselibrary{skins,listings}

\tcbset{
%  skin=enhanced,
%  segmentation style={solid,line width=.4pt},
  boxsep=0pt,
  boxrule=.4pt,
  left=4pt,
  top=3pt,
  right=4pt,
  bottom=2pt,
  colback=black!15,
  after=\global\setlength{\exampleparindent}{0pt}\par\smallskip\noindent,
  }

\lstdefinestyle{tcblatex}{
  literate={ö}{{\"o}}1{ü}{{\"u}}1{ó}{{\'o}}1{ő}{{\H o}}1{ú}{{\'u}}1
  {ű}{{\H u}}1{é}{{\'e}}1{á}{{\'a}}1{í}{{\'i}}1{Ö}{{\"O}}1{Ü}{{\"U}}1
  {Ó}{{\'O}}1{Ő}{{\H O}}1{Ú}{{\'U}}1{Ű}{{\H U}}1{É}{{\'E}}1{Á}{{\'A}}1,
  %{Í}{{\'I}}1,
  comment=[l][\color{gray}]{\%},
  escapeinside={(*}{*)},
  basicstyle=\ttfamily,
  columns=fullflexible,
  keepspaces,
  belowskip=0pt,
  aboveskip=0pt,
  }

\newtcblisting{latexcode}{
  listing only,
  listing options={style=tcblatex},
  }

\newtcblisting{latexexample}{
  sidebyside,
%  sidebyside align=top,
  sidebyside gap=3mm,
  colback=white,
  bicolor,
  colback=black!5,
  colbacklower=white,
  before lower=\setlength{\parindent}{\exampleparindent},
  }

\newlength{\exampleparindent}
\setlength{\exampleparindent}{0pt}
\def\latexexampleparindent{\setlength{\exampleparindent}{10pt}}

\definecolor{hured}{RGB}{207,36,54}
\definecolor{hugreen}{RGB}{69,113,79}
\usepackage{contour}
\def\hulogo{%
  {\contourlength{.8pt}%
  \sbox0{\contour{black!80}{\Huge\bfseries\sffamily\color{black!10}BABEL\lower-3pt\hbox{-}HUNGARIAN}}%
  \mbox{}%
  \lower-\dimeval{2\ht0/3}%
  \rlap{{\color{hured}\rule{\wd0}{\dimeval{\ht0/3}}}}%
  {\color{hugreen}\rule{\wd0}{\dimeval{\ht0/3}}}%
  \llap{\usebox0}}}

\begin{document}

\title{%
   \hulogo\\[1em] 
   Babel support for Hungarian\\
   {\large v1.6c (2026-07-04)}
   }

\author{%
  \textsc{Péter Szabó} and \textsc{Tibor Tómács}\\
  \href{mailto:tomacs.tibor@gmail.com}{\texttt{tomacs.tibor@gmail.com}}\\
  \url{https://github.com/tibortomacs/babel-hungarian}\\
  \url{https://ctan.org/pkg/babel-hungarian}\\[6mm]
  %
  This manual is based on\\%[2mm]
  \emph{Typing Hungarian text with Magyar\LaTeX}%
  \thanks{The original version of this documentation was written in Hungarian in 2009 under the title \emph{``Magyar nyelvű szöveg szedése Magyar\LaTeX-hel''}.}\\
  by \textsc{Péter Szabó}\\
  \url{https://math.bme.hu/latex/magyarldf-doc.pdf}
  }

\date{}

\selectlanguage{english}

\maketitle

\begin{abstract}
This manual documents \textsf{babel} language support for Hungarian as provided by the \textsf{babel-hungarian} package.
The documentation is written in Hungarian, as it is intended primarily for Hungarian-speaking users.
\end{abstract}

\selectlanguage{magyar}

\section{A \textsf{babel} magyar nyelvi támogatásának betöltése}

\subsection{Általános információk}

A \textsf{babel} csomag magyar nyelvi támogatása a \texttt{magyar} vagy \texttt{hungarian} opcióval érhető el:

\begin{latexcode}
\usepackage[magyar]{babel}
\end{latexcode}
vagy

\begin{latexcode}
\usepackage[hungarian]{babel}
\end{latexcode}
Ezek a \textsf{magyar.ldf} fájlt használják a magyar nyelvi beállításokhoz.
A \textsf{magyar.ldf} működése testre szabható opciókkal.
Erre két lehetőség van.
Ha a \textsf{babel} csomag \texttt{magyar} opcióval van betöltve, akkor

\begin{latexcode}
\def\magyarOptions{(*\param{opciók}*)}
\end{latexcode}
vagy

\begin{latexcode}
\PassOptionsToPackage{(*\param{opciók}*)}{magyar.ldf}
\end{latexcode}
illetve \texttt{hungarian} használata esetén

\begin{latexcode}
\def\hungarianOptions{(*\param{opciók}*)}
\end{latexcode}
vagy

\begin{latexcode}
\PassOptionsToPackage{(*\param{opciók}*)}{hungarian.ldf}
\end{latexcode}

Ezeket a parancsokat még a \verb|\usepackage[···]{babel}| előtt ki kell adni!
Az opciók \param{kulcs}\texttt{=}\param{érték} alakúak és több opció esetén azokat vesszővel kell elválasztani.
Az opciók megadási sorrendje lényeges.

Ha több \verb|\def\magyarOptions| (illetve \verb|\def\hungarianOptions|) szerepel a dokumentumban, akkor csak a legutolsó (de még a \verb|\usepackage[···]{babel}| előtti) számít.
Ha több \verb|\PassOptionsToPackage| szerepel, akkor mindegyik számít, futási sorrendben.
Az egy parancson belüli opciólistában a sorrend balról jobbra halad.
Ha egy \param{kulcs} többször szerepel, akkor a sorrend szerinti legutolsó értéke lép érvénybe.
Van néhány olyan \param{kulcs} (pl.~\verb|defaults|), amely más kulcsok értékét módosítja.

Amennyiben a \verb|\PassOptionsToPackage| segítségével vannak megadva az opciók, és valamelyik érték egy parancs, akkor az elé be kell szúrni a \verb|\noexpand| parancsot, különben hibás lesz a fordítás (pl.~\verb|shorthandcs=\noexpand\shu|).

A pdf\LaTeX{} fordító használatakor az ékezetes betűket
tartalmazó szavak elválasztásához be kell tölteni a \textsf{fontenc} csomagot \texttt{T1} opcióval, de Lua\LaTeX{} vagy Xe\LaTeX{} esetén erre nincs szükség.

\subsection{A \textsf{magyar.ldf} ajánlott betöltése}

A \textsf{magyar.ldf} következő betöltése jól illeszkedik a magyar tipográfiai szabályokhoz, ezért a legtöbb esetben ez ajánlott:

\begin{latexcode}
\usepackage[T1]{fontenc} % Csak pdflatex esetén kell!
\PassOptionsToPackage{defaults=hu-min}{magyar.ldf}
\usepackage[magyar]{babel}
\end{latexcode}
Ezzel egyenértékű:

\begin{latexcode}
\usepackage[T1]{fontenc} % Csak pdflatex esetén kell!
\PassOptionsToPackage{defaults=hu-min}{hungarian.ldf}
\usepackage[hungarian]{babel}
\end{latexcode}

\subsection{A \textsf{magyar.ldf} opciói}

Ebben az alszakaszban ismertetésre kerülnek a \textsf{magyar.ldf} opciói.

\begin{description}

\begin{hunopt}{active}{safe}
\item[safe]
Aktívvá teszi az \texttt{activeprefix=} opcióban megadott karaktert, továbbá elérhetővé teszi a \texttt{shorthandcs=} opcióban megadott parancsot.
Az aktív karakterekről \aref{sec-aktiv}.~szakaszban lesz szó részletesen. 
\item[none]
Nem tesz aktívvá karaktert.
Csak \texttt{shorthandcs=none} esetén használható!
\item[onlycs]
Nem aktivál karaktert, de bevezeti a \texttt{shorthandcs=} opcióban megadott nevű parancsot, amely ugyanazt nyújtja hosszabb gépelés árán, mint egy aktívvá tett karakter.
\end{hunopt}

\begin{hunopt}{activeprefix}{babelopt2}
\item[babelopt2]
Az \texttt{active=safe} opció esetén van hatása.
A sima aposztrófot (\verb|'|) teszi aktívvá, ha a \textsf{babel} csomag \texttt{activeacute} opcióval lett betöltve, egyébként pedig a fordított aposztrófot (\verb|`|).
\item[babelopt3]
Az \texttt{active=safe} opció esetén van hatása.
A sima aposztrófot (\verb|'|) teszi aktívvá, ha a \textsf{babel} csomag \texttt{activeacute} opcióval lett betöltve.
A fordított aposztrófot (\verb|`|) teszi aktívvá, ha a \textsf{babel} csomag \texttt{activegrave} opcióval lett betöltve.
Minden más esetben az írógép idézőjelet (\verb|"|) aktiválja.
\item[acute]
Az \texttt{active=safe} opció esetén a sima aposztrófot (\verb|'|) teszi aktívvá.
\item[grave]
Az \texttt{active=safe} opció esetén a fordított aposztrófot (\verb|`|) teszi aktívvá.
\item[quotedbl]
Az \texttt{active=safe} opció esetén az írógép idézőjelet (\verb|"|) teszi aktívvá.
\item[none]
Egyetlen karaktert sem tesz aktívvá.
\end{hunopt}

\begin{hunopt}{activespace}{none}
\item[safe]
Rövid szóköz lesz kihagyva minden kettőspont, pontosvessző, kérdőjel és felkiáltójel előtt.
A magyar tipográfia megköveteli ezt a kihagyást.
Ez az előbbi négy karakter aktiválásával valósul meg.
\item[none]
A kettőspont, pontosvessző, kérdőjel és felkiáltójel előtt nem lesz rövid szóköz kihagyva, ezeket a karakterek nem teszi aktívvá.
\end{hunopt}

\begin{hunopt}{afterindent}{unchanged}
\item[force-yes]
A szakaszcím után az első sort beljebb kezdi a \verb|\parindent| mértékével.
\item[force-no]
A szakaszcím után az első sort nem kezdi beljebb.
\item[unchanged]
Nem történik módosítás.
\end{hunopt}

\begin{hunopt}{amspostsectiondot}{no}
\item[no]
Bizonyos AMS dokumentumosztályok (\textsf{amsproc}, \textsf{amsbook} és \textsf{amsart}) ponttal fejezik be néhány strukturális egység címsorát, amit a magyar tipográfia megtilt.
Ez az opció eltünteti ezt a pontot.
\item[unchanged]
Nem történik módosítás.
\end{hunopt}

\begin{hunopt}{amstocnumlang}{all}
\item[all]
Az \texttt{amstocnumskip=} opció a nem magyar nyelvű fejezetek tartalomjegyzékbeli címsoraira is vonatkozzon.
\item[hu]
Csak a magyar nyelvű fejezetek tartalomjegyzékbeli címsoraira alkalmazza az \texttt{amstocnumskip=} opciót.
\end{hunopt}

\begin{hunopt}{amstocnumskip}{\textbackslash enskip}
\item[\param{hosszúságparancs}]
Előírja, hogy az AMS dokumentumosztályok magyar nyelvű bejegyzéseinek számait a tartalomjegyzékben mekkora kihagyás válassza el a fejezet nevétől.
Tetszőleges hosszúságparancs megadható.
A magyar tipográfia az \verb|\enskip|-et javasolja az AMS által használt \verb|\quad| helyett.
\item[\param{üres}]
Ha az értéket üresen hagyjuk, akkor nem változtat, marad az eredeti \verb|\quad| hosszúság.
\end{hunopt}

\begin{hunopt}{appendixdot}{yes}
\item[yes]
A \verb|\chapter|-rel megadott függelékek sorszáma után rak pontot.
\item[no]
A \verb|\chapter|-rel megadott függelékek sorszáma után nem rak pontot.
\end{hunopt}

\begin{hunopt}{az}{weak}
\item[weak]
A határozott névelőt (a/az) automatikusan kirakni képes \verb|\az| parancsot és az őt használó \verb|\aref| és egyéb parancsokat definiálja, a részleteket lásd \aref{sec-nevelo}.~szakaszban.
Ezeket a parancsokat csak akkor definiálja, ha korábban nem voltak definiálva.
\item[yes]
Ugyanaz mint a \texttt{weak}, de ha korábban ezek a parancsok már definiáltak voltak, akkor azokat átdefiniálja.
\item[no]
Nem definiál határozott névelőre vonatkozó parancsokat.
\end{hunopt}

\begin{hunopt}{captions}{hu}
\item[hu]
A járulékos fejezetneveket és egyéb feliratokat magyarul jeleníti meg.
\item[unchanged]
Nem történik módosítás.
\end{hunopt}

\begin{hunopt}{chapterhead}{yes}
\item[yes]
A fejezetek elején a magyarnak megfelelő formátumban először a sorszám, egy pont, majd a ,,fejezet'' szó jelenik meg.
Az alapértelmezett angol formátumban ez fordítva van pont nélkül.
\item[unchanged]
Nem történik módosítás.
\end{hunopt}

\begin{hunopt}{chapternumber}{unchanged}
\item[Huordinal]
A fejezetszámok magyarul betűzve jelennek meg, az elején nagybetűvel, például ,,Első fejezet''.
\item[huordinal]
A fejezetszámok magyarul betűzve jelennek meg, az elején kisbetűvel, például ,,első fejezet''.
\item[unchanged]
Nem történik módosítás.
\end{hunopt}

\begin{hunopt}{classmod}{yes}
\item[yes]
A \textsf{book}, a \textsf{report} és az \textsf{article} osztályok viselkedését képes módosítani.
Ezekben a \verb|\part| és \verb|\chapter| parancsok, a tartalomjegyzékek és a fejlécek tipográfiáját módosítja.
Bár ismeri a \textsf{letter} osztályt is, a levél fejlécét és egyéb formai elemeit nem a magyar szabályok szerint helyezi el, csak a címsort módosítja.
\item[unchanged]
Nem történik módosítás.
\end{hunopt}

\begin{hunopt}{dottedtocline}{fix}
\item[fix]
A \verb|\@dottedtocline| parancs átdefiniálásával engedélyezi, hogy a tartalomjegyzékek pontozott sorainak végén az oldalszám balra túllógjon a számára fenntartott \verb|\@pnumwidth| távolságon.
Az opció nem kötődik szorosan a magyar nyelvhez.
\item[unchanged]
Nem történik módosítás.
\end{hunopt}

\begin{hunopt}{emitdate}{weak}
\item[weak]
A dátumok magyar helyesírás szerinti megjelenítését végző, több formából választani engedő és a ragozást is támogató \verb|\emitdate| parancsot definiálja, feltéve, hogy korábban nem volt definiálva.
A részleteket lásd \aref{sec-datum}.~szakaszban.
Lásd még a \texttt{hutoday=} opciót.
\item[yes]
Ugyanaz mint a \texttt{weak}, de ha korábban az \verb|\emitdate| már definiált volt, akkor azt átdefiniálja.
\item[no]
Nem definiálja az \verb|\emitdate| parancsot.
\end{hunopt}

\begin{hunopt}{extras}{yes}
\item[yes]
A magyar nyelv aktívvá válásakor számos beállítás módosul az \verb|\extrasmagyar| makró futtatásával, továbbá amikor egy másik nyelv lesz aktív, akkor előbb a \verb|\noextrasmagyar| makró fut le.
A \textsf{magyar.ldf} betöltődésekor a fenti két makróba építi bele saját nyelvspecifikus beállításait.
Az opció \texttt{yes} értéke engedi érvényre jutni ezeket a beállításokat
\item[no]
Letiltja az \verb|\extrasmagyar| és \verb|\noextrasmagyar| makrók futását.
Ez egyedül hibakeresési céllal javasolt, mert épp a \textsf{magyar.ldf} fő célját hiúsítja meg.
Hibakeresés esetén érdemes az \texttt{extras=no}-t is tartalmazó \texttt{defaults=safest} alapbeállításokból kiindulni (lásd \aref{subsec-defaults}.~alszakaszt).
\end{hunopt}

\begin{hunopt}{fancyhdr}{hu}
\item[hu]
A \textsf{fancyhdr} csomag viselkedését igazítja a magyar tipográfiához.
\item[unchanged]
Nem történik módosítás.
\end{hunopt}

\begin{hunopt}{figurecaptions}{hu}
\item[hu]
Az ábra feliratokat igazítja a magyar tipográfiához.
\item[unchanged]
Nem történik módosítás.
\end{hunopt}

\begin{hunopt}{footnote}{yes}
\item[yes]
Definiálja a \verb|\footnotestyle| parancsot, mellyel \aref{sec-labjegyzet}.~szakaszban leírt módon testre szabható a lábjegyzetek megjelenítése.
\item[huplain]
Ugyanaz, mint a \texttt{yes}, de ezen felül még a magyar tipográfiának megfelelő, arab számokat használó stílust állít be.
\item[starplain]
Ugyanaz, mint a \texttt{yes}, de csillagokat használó stílust állít be, ami szintén megfelel a magyar tipográfiának.
\item[unchanged]
Nem változtat a lábjegyzeteken, és a \verb|\footnotestyle| parancsot sem definiálja.
\end{hunopt}

\begin{hunopt}{frenchspacing}{unchanged}
\item[yes]
Az írásjelek utáni szóköz szélességét a magyar tipográfiának megfelelően egyenletesre állítja (\verb|\frenchspacing|), amíg a magyar nyelv aktív.
\item[no]
A magyar nyelv aktívvá válásakor az írásjeltől függő változó szóközszélességet ír elő. (Nem felel meg a magyar tipográfiának!)
\item[unchanged]
Nem történik módosítás.
\end{hunopt}

\begin{hunopt}{hang}{weak}
\item[weak]
A szimbólummal kezdődő, függő bekezdések nyitására szolgáló \verb|\hang| parancsot definiálja, amennyiben korábban nem volt definiálva.
A részleteket lásd \aref{sec-idezet-parbeszed}.~szakaszban.
\item[yes]
Ugyanaz mint a \texttt{weak}, de ha korábban a \verb|\hang| már definiált volt, akkor azt átdefiniálja.
\item[no]
Nem definiálja a \verb|\hang| parancsot.
\end{hunopt}

\begin{hunopt}{hunnewlabel}{yes}
\item[yes]
A generált oldalszámok és fejezetszámok névelővel történő ellátását segítő \verb|\hunnewlabel| segédmakrót definiálja.
A részleteket lásd \aref{sec-nevelo}.~szakaszban.
\item[no]
Nem generál \verb|\hunnewlabel|-eket.
\item[only-hu]
Csak akkor generál \verb|\hunnewlabel|-eket, ha a \verb|\label| parancs hívásakor a magyar nyelv aktív.
\end{hunopt}

\begin{hunopt}{hunumbers}{yes}
\item[yes]
Az egész számokat $-9999\dots9999$ között magyarul kibetűzni képes parancsokat definiálja: \verb|\@hunumeral|, \verb|\@Hunumeral|, \verb|\@huordinal|, \verb|\@huordinal|.
Ezek a \texttt{partnumber=} és \texttt{chapternumber=} opciók \texttt{huordinal} és \texttt{Huordinal} értékeihez is kellenek, de önállóan is használhatóak.
Például
\begin{latexexample}
\makeatletter
\@hunumeral{123}\\
\@Hunumeral{123}\\
\@huordinal{123}\\
\@Huordinal{123}
\makeatother
\end{latexexample}
Számok kibetűzéséhez a \textsf{numspell} csomag (\url{https://ctan.org/pkg/numspell}) is használható, amely a magyaron kívül több nyelven is működik.
\item[no]
Nem definiálja az előbbi parancsokat.
\end{hunopt}

\begin{hunopt}{hutoday}{yes}
\item[yes]
A dátumok magyar megjelenítéséért felelős parancsokat definiálja: \verb|\today|, \verb|\ondatemagyar|, \verb|\ontoday|.
Bővebb leírás olvasható \aref{sec-datum}.~szakaszban.
A magyar nyelvű dátummegjelenítésnek egy gazdagabb, ragozást is megengedő formája az \verb|\emitdate| parancs, melyet az \texttt{emitdate=} opcióval lehet engedélyezni.
\item[no]
Nem módosít az alapbeállításokon.
\end{hunopt}

\begin{hunopt}{hynumberline}{unchanged}
\item[hu]
Ha a \textsf{hyperref} csomagot \texttt{bookmarksnumbered} opcióval töltötte be, akkor a PDF könyvjelzőjében a szakasz számai után pont fog állni a magyar tipográfiának megfelelően.
\item[unchanged]
Nem történik módosítás.
\end{hunopt}

\begin{hunopt}{hyphenmins}{22}
\item[\param{számjegy bal\,}\param{számjegy jobb}]
A \param{számjegy bal\,} illetve \param{számjegy jobb} azt adja meg, hogy a magyar nyelvű szavak bal illetve jobb oldalán legalább hány karakter közé nem szabad automatikusan elválasztójelet tenni.
A \textsf{babel} és a \LaTeX\ alapbeállítása 23, a \textsf{magyar.ldf}-é pedig 22, tehát például ,,Aliz'' a bal oldala miatt nem választható el, ,,Sára'' viszont igen, mivel mindkét oldalon megvan a két karakter.
\item[unchanged]
Nem történik módosítás.
\end{hunopt}

\begin{hunopt}{labelenums}{unchanged}
[Az \texttt{enumerate} környezettel létrehozott lista szintjeinek számozását állítja be.]
\item[hu-a]
A szintek száma 5.
Ezek számozási típusa sorrendben:
I. 1. \textit{a)} \textit{$\alpha$)} \textit{A)}
\item[hu-A]
A szintek száma 5.
Ezek számozási típusa sorrendben:
I. 1. \textit{A)} \textit{a)} \textit{$\alpha$)}
\item[hu-d]
A szintek száma 4.
Ezek számozási típusa sorrendben:
1. \textit{a)} \textit{$\alpha$)} \textit{A)}
\item[unchanged]
Nem történik módosítás.
\end{hunopt}

\begin{hunopt}{labelitems}{unchanged}
[Az \texttt{itemize} környezettel létrehozott lista szintjeinek jelöléseit állítja be.]
\item[hu]
A szintek száma 5.
Ezek jelei sorrendben:
\labelitemi\ \labelitemii\ \labelitemiii\ \labelitemiv\ \labelitemv
\item[unchanged]
Nem történik módosítás.
\end{hunopt}

\begin{hunopt}{longcaption}{justified}
[Az úsztató környezetek feliratának formátumát szabályozza.
Nincs hatása, amennyiben a \textsf{caption} vagy \textsf{caption2} csomagok valamelyikét betölti.]
\item[justified]
A címke végére kettőspont helyett pontot rak.
A többsoros címek sorkizártan jelennek meg.
\item[centered]
A címke végére kettőspont helyett pontot rak.
A többsoros címek középre zártan jelennek meg.
\item[centernewline]
Annyiban különbözik a \texttt{centered} értéktől, hogy ekkor a címke után sortörés van.
\item[unchanged]
Nem történik módosítás.
\end{hunopt}

\begin{hunopt}{mathbrk}{define}
\item[define]
A matematikai szimbólumok sortörését támogató \verb|\MathBrk| parancsot definiálja.
A részleteket lásd \aref{sec-matek}.~szakaszban.
\item[fix]
A matematikai szimbólumok sortörését támogató \verb|\MathBrk| parancsot definiálja, és ezen felül még a magyar tipográfiai szokásokhoz igazítja néhány matematikai
bináris operátor és bináris reláció sortöréssel kapcsolatos viselkedését.
A részleteket lásd \aref{sec-matek}.~szakaszban.
\item[unchanged]
Nem történik módosítás.
\end{hunopt}

\begin{hunopt}{mathfactorial}{define}
\item[define]
A matematikai faktoriális jelet helyes térközzel beszúró \verb|\factorial| parancsot definiálja.
A részleteket lásd \aref{sec-matek}.~szakaszban.
\item[fix]
A matematikai faktoriális jelet helyes térközzel beszúró \verb|\factorial| parancsot definiálja és ezen felül még átírja a felkiáltójel képleten belüli jelentését \verb|\factorial|-ra.
Ha az \texttt{activespace=} és a \texttt{mathfactorial=} opciókat együtt használjuk, akkor a felkiáltójelre matematikai képletben a \texttt{mathfactorial=}, szövegben pedig az \texttt{activespace=} van hatással.
A részleteket lásd \aref{sec-matek}.~szakaszban.
\item[unchanged]
Nem történik módosítás.
\end{hunopt}

\begin{hunopt}{mathhucomma}{define}
[A vessző viselkedését szabályozza matematikai üzemmódban.]
\item[fix]
A magyar tipográfia a tizedes törtekben tizedespont helyett tizedesvesszőt használ. Magyar tizedes törtek alapesetben nem gépelhetők közvetlenül \verb|$1,2$| módon, mert a vessző körül kihagyott térköz hibás lenne (ugyanis a vessző operátorosztálya Punct, a jó osztály, ahová a pont is tartozik, pedig Ord).

Ezzel az opcióval azonban, amennyiben egy vesszőt közvetlenül szám követ matematikai üzemmódban, akkor a vessző Ord osztályba sorolódik, így az \verb|$1,2$| már valóban tizedes törtet fog jelölni.

Ha a vessző nem tizedesvesszőt jelent és szám követi, akkor be kell tenni egy szóközt a vessző és szám közé.
Pl.~\verb|$1, 2, 3$|.
Ekkor a vessző az eredeti Punct jelként működik.

Ha a vesszőt közvetlenül nem szám, hanem például egy betű követi, akkor marad a vessző eredeti Punct jelként való besorolása.
Pl.~\verb|$a,b,c$|.

Az előbb leírtak úgy valósulnak meg, hogy definiálódik egy \verb|\HuComma| parancs, ami a magyar tizedesvessző szerepét látja el, és ezután a vessző képleten belüli jelentését átírja \verb|\HuComma|-ra.
\item[define]
A \texttt{fix} értéknél említett \verb|\HuComma| definiálódik, de a vessző képleten belüli jelentését nem írja át.
\item[unchanged]
Nem történik módosítás.
Ebben az esetben pl.~\verb|$1{,}2$| módon lehet magyar tizedes törtet gépelni.
\end{hunopt}

\begin{hunopt}{mathmuskips}{unchanged}
\item[hu]
A \verb|\thickmuskip|, \verb|\medmuskip| és \verb|\thinmuskip| távolságok állításával igyekszik megfelelni a magyar tipográfiának -- bár teljesen nem tud, mert a magyar tipográfia a \TeX{} térköz-generálásától lényegesen eltérően építi fel a szabályokat, melyek követése ráadásul kevésbé tetszetős képleteket eredményezne, mint ez az opció.
\item[latex]
Az előbbi távolságokat a \LaTeX{} alapbeállításaira állítja.
\item[unchanged]
Nem történik módosítás.
\end{hunopt}

\begin{hunopt}{mathreal}{weak}
\item[weak]
A képletekben a tizedespontot tizedesvesszővel helyettesítő \verb|\MathReal| parancsot definiálja, amennyiben korábban nem volt definiálva.
A részleteket lásd \aref{sec-matek}.~szakaszban.
\item[yes]
Ugyanaz mint a \texttt{weak}, de ha korábban a \verb|\MathReal| már definiált volt, akkor azt átdefiniálja.
\item[no]
Nem definiálja a \verb|\MathReal| parancsot.
\end{hunopt}

\begin{hunopt}{mond}{weak}
\item[weak]
Párbeszédekben a megszólalást bevezető \verb|\mond| parancsot definiálja, amennyiben korábban nem volt definiálva.
A részleteket lásd \aref{sec-idezet-parbeszed}.~szakaszban.
\item[yes]
Ugyanaz mint a \texttt{weak}, de ha korábban a \verb|\mond| már definiált volt, akkor azt átdefiniálja.
\item[no]
Nem definiálja a \verb|\mond| parancsot.
\end{hunopt}

\begin{hunopt}{openqq}{maybedown}
[Azt adja meg, hogy aktív karakter után írt fordított aposztróf (\texttt{\textasciigrave}) illetve sima aposztróf (\texttt{\textquotesingle}) hogyan viselkedjen.
Például, ha a fordított aposztróf aktív karakter, akkor mi legyen a \texttt{\textasciigrave\textasciigrave} illetve \texttt{\textasciigrave\textquotesingle} eredménye a PDF-ben?]
\item[down]
Ezzel nem használható az \texttt{active=none} opció! Hatása:\\
\param{aktív karakter}+\verb|`| $=$ \verb|\quotedblbase| (\quotedblbase)\\
\param{aktív karakter}+\verb|'| $=$ \verb|\textquotedblleft| (\textquotedblleft)
\item[up]
Ezzel nem használható az \texttt{active=none} opció! Hatása:\\
\param{aktív karakter}+\verb|`| $=$ \verb|\textquotedblleft| (\textquotedblleft)\\
\param{aktív karakter}+\verb|'| $=$ \verb|\quotedblbase| (\quotedblbase)
\item[unchanged]
Nem kezeli parancsként az \param{aktív karakter}+\verb|`| és \param{aktív karakter}+\verb|'| párosításokat.
\item[maybedown]
Az \texttt{active=none} opció esetén az \texttt{unchanged}, egyéb esetben pedig a \texttt{down} értékkel ekvivalens.
\end{hunopt}

\begin{hunopt}{partnumber}{unchanged}
\item[Huordinal]
A részek számai magyarul betűzve jelennek meg, az elején nagybetűvel, például ,,Első rész''.
\item[huordinal]
A részek számai magyarul betűzve jelennek meg, az elején kisbetűvel, például ,,első rész''.
\item[unchanged]
Nem történik módosítás.
\end{hunopt}

\begin{hunopt}{postdescription}{unchanged}
[A \texttt{description} leíró lista környezetben az \texttt{\textbackslash item} opcionális paraméterével megadott listaelem címe és az ezt követő szöveg közti kihagyást, illetve a díszjelet állítja be.]
\item[dot]
A díszjel egy pont, ami a listaelem címének betűtípusával jelenik meg.
A pont után 0,5\,em térköz van.
\item[mddot]
Ugyanaz, mint a \texttt{dot}, de mindig normál vastagságú fonttal.
\item[bfdot]
Ugyanaz, mint a \texttt{dot}, de mindig félkövér vastagságú fonttal.
\item[diamond]
A díszjel \mbox{\raise.5ex\hbox{\font\f=msam10 at1.3ex \f\char7}}, ami a listaelem címének betűtípusával jelenik meg.
Előtte és utána 0,5\,em térköz van.
\item[star]
A díszjel egy csillag.
Előtte és utána 0,5\,em térköz van.
\item[endash]
A díszjel egy nagykötőjel, ami a listaelem címének betűtípusával jelenik meg.
Előtte és utána normál szóköz van.
\item[colon]
A díszjel egy kettőspont, ami a listaelem címének betűtípusával jelenik meg.
Ezután 0,5\,em térköz van.
\item[quad]
Nincs díszjel, csak 1\,em térköz.
\item[enskip]
Nincs díszjel, csak 0,5\,em térköz.
\item[unchanged]
Nem történik módosítás.
\end{hunopt}

\begin{hunopt*}{postpara}{unchanged}
Abban különbözik a \texttt{postdescription=} opciótól, hogy a \texttt{\textbackslash paragraph} címsorára van hatással.
A lehetséges értékek megegyeznek a \texttt{postdescription=} értékeivel.
\end{hunopt*}

\begin{hunopt*}{postsubpara}{unchanged}
Abban különbözik a \texttt{postdescription=} opciótól, hogy a \texttt{\textbackslash subparagraph} címsorára van hatással.
A lehetséges értékek megegyeznek a \texttt{postdescription=} értékeivel.
\end{hunopt*}

\begin{hunopt}{refstruc}{weak}
\item[weak]
A dokumentum strukturális egységeire hivatkozó \verb|\refstruc| parancsot definiálja, feltéve, hogy korábban nem volt definiálva.
A részleteket lásd \aref{sec-hivatkozas}.~szakaszban.
\item[yes]
Ugyanaz mint a \texttt{weak}, de ha korábban a \verb|\refstruc| már definiált volt, akkor azt átdefiniálja.
\item[no]
Nem definiálja a \verb|\refstruc| parancsot.
\end{hunopt}

\begin{hunopt}{sectiondot}{safe}
\item[safe]
Pontot tesz az ábra/táblázat számok után az ábra/táblázat jegyzékekben, illetve a fejezet/szakasz számok után a tartalomjegyzékben.
\item[none]
Nem történik módosítás.
\end{hunopt}

\begin{hunopt}{shorthandcs}{\textbackslash shu}
\item[\param{parancs}]
Az \texttt{active=none} opcióval nem használható együtt.
A megadott \param{parancs} az aktív karakterrel lesz egyenértékű.
Fontos, hogy ekkor a \param{parancs} elnyeli az első paraméterét.
Például \verb|shorthandcs=\shu| esetén \verb|\shu&tty| vagy \verb|\shu`tty| hatása ugyanaz, mint \verb|`tty| (feltéve, hogy \verb|`| aktív karakter).
\item[\param{parancsnév}]
Ugyanaz, mint az előbbi, de itt csak a parancs nevét kell megadni \verb|\| nélkül.
Ebben az esetben a definiált parancs nem nyeli el az első paraméterét.
Például \verb|shorthandcs=shu| esetén \verb|\shu tty| hatása ugyanaz, mint \verb|`tty| (feltéve, hogy \verb|`| aktív karakter).
\item[none]
Nem lesz ilyen parancs.
\end{hunopt}

\begin{hunopt}{shortrefcmds}{yes}
\item[yes]
Háromkarakteres neveket definiál a névelővel hivatkozó parancsoknak: \verb|\aref| $\to$ \verb|\azr|, \verb|\Aref| $\to$ \verb|\Azr|, \verb|\apageref|
$\to$ \verb|\azp|, \verb|\Apageref| $\to$ \verb|\Azp|, \verb|\acite| $\to$ \verb|\azc|, \verb|\Acite| $\to$ \verb|\Azc|.
A határozott névelőről szóló \ref{sec-nevelo}.~szakaszban további részletek olvashatók.
\item[no]
Ezeket a hárombetűs parancsokat nem definiálja.
\end{hunopt}

\begin{hunopt}{suggestions}{yes}
\item[yes]
Figyelmeztető üzeneteket jelenít meg a naplófájlban, ha a \textsf{magyar.ldf}-et hiányosan vagy hibásan próbálják használni.
\item[no]
Nincs ilyen figyelmeztetés.
\end{hunopt}

\begin{hunopt}{tablecaptions}{hu}
\item[hu]
A táblázatok feliratait igazítja a magyar tipográfiához.
\item[unchanged]
Nem történik módosítás.
\end{hunopt}

\begin{hunopt}{textqq}{weak}
\item[weak]
Az idézőjeleket generáló \verb|\textqq| parancsot definiálja, feltéve, hogy korábban nem volt definiálva.
A részleteket \aref{sec-idezet-parbeszed}.~szakasz tartalmazza.
\item[yes]
Ugyanaz mint a \texttt{weak}, de ha korábban a \verb|\textqq| már definiált volt, akkor azt átdefiniálja.
\item[no]
Nem definiálja a \verb|\textqq| parancsot.
\end{hunopt}

\begin{hunopt}{theoremtitle}{hu}
\item[hu]
A magyar tipográfiai szabályoknak megfelelően szedi a tételszerű környezetek címeit: a sorszám után pont és szóköz van, majd a környezet címe (tétel, lemma,
bizonyítás, stb.), megint egy pont, és végül szóköz jön.
Ez az opció kompatibilis a \textsf{theorem}, \textsf{ntheorem} és \textsf{amsthm} csomagokkal.
A \textsf{theorem} és \textsf{ntheorem} esetén a magyarítást a \verb|\theoremstyle{magyar-plain}| illetve \verb|\theoremstyle{hungarian-plain}| parancs végzi attól függően, hogy a \textsf{babel} csomag \texttt{magyar} vagy \texttt{hungarian} opcióval lett betöltve.
Ez a parancs a dokumentumtest elején aktiválódik, így a \verb|\begin{document}| után írt \verb|\newtheorem| esetén működik.
\item[unchanged]
Nem történik módosítás.
\end{hunopt}

\begin{hunopt}{titles}{\textbackslash enskip}
\item[\param{hosszúságparancs}]
A strukturális egységek címeinek (pl.~\verb|\section|) szedésekor a pont és a cím közötti térközt adja meg.
\item[unchanged]
Nem történik módosítás.
\end{hunopt}

\begin{hunopt}{told}{weak}
\item[weak]
Az egész számok toldalékolását végző \verb|\told| parancsot definiálja, feltéve, hogy korábban nem volt definiálva.
A részleteket \aref{sec-toldalek}.~szakasz tartalmazza.
\item[yes]
Ugyanaz mint a \texttt{weak}, de ha korábban a \verb|\told| már definiált volt, akkor azt átdefiniálja.
\item[no]
Nem definiálja a \verb|\told| parancsot.
A \texttt{varioref} csomag magyarítása használja a \verb|\told| parancsot, ezért ilyenkor a \texttt{told=no} nem megfelelő választás.
\end{hunopt}

\end{description}

\subsection{Az opciók alapértékeinek beállítása}\label{subsec-defaults}

A következő opció \aref{tablazat-defaults}.~táblázatnak megfelelően állítja be a korábban ismertetett opciók alapértékeit.

\begin{description}

\begin{hunopt}{defaults}{over-1.4}
\item[over-1.4]
A \textsf{magyar.ldf} 1.4-es verziójának tipográfiájától csak kevéssé tér el, de néhány fontos ponton van különbség.
Az új, más csomagokkal össze nem akadó parancsokat elérhetővé teszi, de nem kapcsolja be.
Célja, hogy az összes 1.4-es verzióhoz készült dokumentum leforduljon vele, és csak kevés, a sor- és oldalhatárokat általában nem érintő látható eltérés legyen az 1.4-es verzióhoz képest.
Ha ehelyett a szerző célja az, hogy a dokumentum a magyar tipográfiai szabályokat minél inkább kövesse, akkor helyette a \texttt{defaults=hu-min} ajánlott.
\item[compat-1.4]
Nem változtat lehetőség szerint a \textsf{magyar.ldf} 1.4-es verziójának tipográfiáján, de kijavít számos megvalósítási hibát és más \LaTeX{} csomagokkal való összeférhetetlenséget.
Ezt ajánljuk azoknak, akik egy régi dokumentumot szeretnének vizuális változás nélkül az új \textsf{magyar.ldf}-fel újrafordítani. 
A teljes egyezés nem garantált, a kompatibilitás a \textsf{magyar.ldf}-nek csak az opciók által befolyásolt részeire terjed ki.
Elrejti az 1.4-es verzió után definiált parancsokat is.
\item[prettiest]
Minden új szolgáltatást bekapcsol, a magyar tipográfiához a lehető legjobban illeszkedik.
Vigyázat!
Az itt bekapcsolt szolgáltatások esetleg összeakadhatnak más csomagokkal.
\item[hu-min]
Az alapbeállításoktól annyiban tér el, hogy nem az 1.4-es verzióval való kompatibilitásra figyel, hanem az opciókat úgy állítja be, hogy azok ne mondjanak ellent a magyar tipográfiai alapelveknek.
Ha a szabályok választást tesznek lehetővé, akkor ez az alapbeállítás a legegyszerűbben megvalósítható és legkevesebb kompatibilitási problémát előhozó értékek közül válogat.
A legtöbb magyar nyelvű dokumentum esetén ez a beállítás javasolt.
A \texttt{defaults=hu-min} által aktivált opciók listáját kiírathatja a fordítási naplóba a \verb|\magyarDumpHuMin| paranccsal.
\item[safest]
Szinte mindent kikapcsol, hatása olyan, mintha a \textsf{magyar.ldf}-et be se töltöttük volna.
Hibakeresésnél érdemes ezt választani.
Váratlan hibaüzenetnél, vagy ha az a gyanú, hogy a \textsf{magyar.ldf} összeakad valamely másik csomaggal, akkor át kell váltani \texttt{defaults=safest}-re, és ha a probléma továbbra is fennáll, akkor nem a \textsf{magyar.ldf}-fel volt baj.
\end{hunopt}

\begin{hunopt}{safest}{no}
\item[no]
Nincs hatása.
\item[yes]
Ugyanaz, mint a \texttt{defaults=safest}.
\end{hunopt}

\end{description}

{\ttfamily\footnotesize 
\begin{longtable}{|l|l|l|l|l|l|}
\caption{Az opciók alapértékei a \texttt{defaults=} opció különböző értékei esetén}
\label{tablazat-defaults}\\
\hline
\textbf{\small\rule{0pt}{.9em}defaults=} & \textbf{\small over-1.4} & \textbf{\small compat-1.4} & \textbf{\small prettiest} & \textbf{\small safest} & \textbf{\small hu-min} \\
\hline
\endfirsthead
\hline
\textbf{\small\rule{0pt}{.9em}defaults=} & \textbf{\small over-1.4} & \textbf{\small compat-1.4} & \textbf{\small prettiest} & \textbf{\small safest} & \textbf{\small hu-min} \\
\hline
\endhead
\hline
\multicolumn{6}{r@{}}{\rule{0pt}{1em}\rmfamily\em Folytatás a következő oldalon!}
\endfoot
\hline
\endlastfoot
%
active= & safe & safe & safe & none & safe \\
activeprefix= & babelopt2 & grave & babelopt2 & none & babelopt2 \\
activespace= & none & none & safe & none & safe \\
afterindent= & unchanged & unchanged & unchanged & unchanged & unchanged \\
amspostsectiondot= & no & unchanged & no & unchanged & no \\
amstocnumlang= & all & all & all & all & all \\
amstocnumskip= & \textbackslash enskip &  & \textbackslash enskip &  & \textbackslash enskip \\
appendixdot= & yes & yes & no & yes & no \\
az= & weak & yes & yes & no & weak \\
captions= & hu & hu & hu & unchanged & hu \\
chapterhead= & yes & yes & yes & no & yes \\
chapternumber= & unchanged & unchanged & unchanged & unchanged & unchanged \\
classmod= & yes & yes & yes & unchanged & yes \\
dottedtocline= & fix & unchanged & fix & unchanged & fix \\
emitdate= & weak & no & yes & no & weak \\
extras= & yes & yes & yes & no & yes \\
fancyhdr= & hu & unchanged & hu & unchanged & hu \\
figurecaptions= & hu & hu & hu & unchanged & hu \\
footnote= & yes & unchanged & yes & unchanged & huplain \\
frenchspacing= & unchanged & unchanged & yes & unchanged & yes \\
hang= & weak & no & yes & no & weak \\
hunnewlabel= & yes & yes & yes & no & yes \\
hunumbers= & yes & no & yes & no & yes \\
hutoday= & yes & yes & yes & no & yes \\
hynumberline= & unchanged & unchanged & hu & unchanged & hu \\
hyphenmins= & 22 & unchanged & 22 & unchanged & 22 \\
labelenums= & unchanged & unchanged & hu-a & unchanged & hu-d \\
labelitems= & unchanged & unchanged & hu & unchanged & hu \\
longcaption= & justified & justified & centered & unchanged & centered \\
mathbrk= & define & unchanged & fix & unchanged & fix \\
mathfactorial= & define & unchanged & fix & unchanged & define \\
mathhucomma= & define & unchanged & fix & unchanged & fix \\
mathmuskips= & unchanged & unchanged & hu & unchanged & unchanged \\
mathreal= & weak & no & yes & no & weak \\
mond= & weak & no & yes & no & weak \\
openqq= & maybedown & maybedown & maybedown & unchanged & maybedown \\
partnumber= & unchanged & unchanged & Huordinal & unchanged & unchanged \\
postdescription= & unchanged & unchanged & dot & unchanged & dot \\
postpara= & unchanged & unchanged & diamond & unchanged & diamond \\
postsubpara= & unchanged & unchanged & star & unchanged & star \\
refstruc= & weak & no & yes & no & weak \\
sectiondot= & safe & safe & safe & none & safe \\
shorthandcs= & \textbackslash shu & none & \textbackslash shu & none & \textbackslash shu \\
shortrefcmds= & yes & yes & no & no & yes \\
suggestions= & yes & no & yes & yes & yes \\
tablecaptions= & hu & hu & hu & unchanged & hu \\
textqq= & weak & no & yes & no & weak \\
theoremtitle= & hu & hu & hu & unchanged & hu \\
titles= & \textbackslash enskip & \quad & \textbackslash enskip & unchanged & \textbackslash enskip \\
told= & weak & no & yes & no & weak \\
\hline
\end{longtable}}

\section{Idézetek, párbeszédek}\label{sec-idezet-parbeszed}

\begin{latexcode}
\textqq{(*\param{szövegközi idézet}*)}
\end{latexcode}
Magyar nyelvű dokumentumokban a következő módon lehet szövegközi idézeteket írni:

\begin{latexexample}
,,Külső >>középső 'belső' 
középső<< külső''
\end{latexexample}
A \texttt{textqq=yes} opció esetén definiálódik a \verb|\textqq| parancs, amivel ez automatizálható:

\begin{latexexample}
\textqq{Külső \textqq{középső 
\textqq{belső} középső} külső}
\end{latexexample}
Magyartól eltérő nyelvre váltásnál a \verb|\textqq| angol tipográfia szerint fog működni, függetlenül a nyelvtől.

\begin{latexcode}
\mond (*\param{kimondott szöveg}*)
\end{latexcode}
Ha egy szereplő a műben megszólal, az általa mondott szöveget új bekezdésben, gondolatjellel kell kezdeni.
A gondolatjel után nem nyúló szóközt kell hagyni.
Ezt valósítja meg \texttt{mond=yes} opció esetén a \verb|\mond| parancs.
A kimondott szövegbe gondolatjelek közt leírást is lehet ékelni, melyet ponttal kell lezárni.
A kimondott szöveg végére szükség esetén ki kell tenni a kérdőjelet vagy a felkiáltójelet, de a pontot tilos.
Például

\latexexampleparindent
\begin{latexexample}
Egy deszkán találta magát,
amely a tenger hullámain
zötykölődött.
\mond Hol a Titanic? -- kérdezte,
de nem kapott választ.
\mond Ez nem lehet -- szólalt
meg ismét. -- Öt perce még a
kabinomban voltam.
\end{latexexample}

\begin{latexcode}
\hang{(*\param{bekezdéskezdő jel}*)}(*\param{bekezdés}*)
\end{latexcode}
Ad hoc szervezésű függő felsorolásokhoz lehet hasznos a \texttt{hang=yes} opció esetén definiált \verb|\hang| parancs, amely az új bekezdés első sorát a bal margón kezdi, a további sorokat pedig a megadott jel mögött.
A \verb|\hang| listakörnyezeteken kívül használandó, ebben is különbözik az \verb|\item|-től.
Például

\begin{latexexample}
\hang{-- }Első sor\\
második sor.\par
\hang{$\bullet\circ$ }Egy másik.
\end{latexexample}

\section{Aktív karakter}\label{sec-aktiv}

Alapértelmezésben a \textsf{magyar.ldf} a fordított aposztróf (\verb|`|) karaktert teszi aktívvá (ez megváltoztatható az \texttt{activeprefix=} opcióval), továbbá definiálja a \verb|\shu| parancsot, és ezután az \verb|\shu`| és a \verb|`| ugyanarra használható.

Az aktív karakter az őt követő karaktert makróparaméterként veszi, tehát például \verb|`=| írható így: \verb|`{=}| vagy így is: \verb*|` =|. 
A \verb|\string| az őt követő aktív karakter speciális jelentését megszünteti.
Így például \verb|\string`\string`| egy angol nyitó idézőjelet eredményez, feltéve, hogy az aktuális fontban van ilyen ligatúra.

A \textsf{magyar.ldf} aktív karaktere az alábbi funkciókat nyújtja:
\begin{longtable}[l]{l|l}
\verb|``|   & ,, magyar nyitó idézőjel,\texttt{ openqq=} miatt eltérhet\\
\verb|`'|   & `' angol nyitó idézőjel,\texttt{ openqq=} miatt eltérhet\\
\verb|`=|   & szón belüli kötőjel, ha a részszavakat is el akarjuk választani\\
\verb|`_|   & szón belüli feltételes kötőjel\\
            & \verb|\showhyphens{bohém\-életre bohém`_életre}| $\to$\\
            & \texttt{bohém-életre bo-hém-élet-re}\\
\verb|`--|  & gondolatjel kis szóközzel, személynevek elválasztására\\
            & \verb|Kiss Előd`--Nagy Pál| $\to$ Kiss Előd`--Nagy Pál\\
\verb|`-|   & ugyanaz, mint \verb|`=|, kivéve matematikai képletben, vagy ha így folytatódik: \verb|`--|\\
            & matematikai képletben törhető szóköz: \verb|$1<2,`-3<4$| $\to$ $1<2,\ 3<4$\\
\verb!`|!   & fontos kötőjel, sor végén és elején is megjelenik\\
\verb|`<|   & `< francia nyitó idézőjel\\
\verb|`>|   & `> francia nyitó idézőjel\\
\verb|`"|   & szón belüli \verb|\allowbreak|, résszavak elválasztása automatikus\\
            & \verb|\showhyphens{kutatás`"+`"fejlesztés}| $\to$ \texttt{ku-ta-tás+fej-lesz-tés}\\
\verb|`~|   & ligatúrát nem képező kötőjel, pl.~\verb|`~`~`~| $\to$ `~`~`~\\
\verb|`ccs| & sor végén ,,ccs'', elválasztva ,,cs-cs''\\
            & \verb|\showhyphens{loccsan szétlo`ccsantja}| $\to$\\
            & \texttt{loccsan szét-locs-csant-ja}\\
\verb|`CCS| & sor végén ,,CCS'', elválasztva ,,CS-CS''\\
\verb|`cxy| & sor végén ,,cxy'', elválasztva ,,cs-xy''; de csak akkor hasznos, ha \texttt{xy = cs}
\end{longtable}
\noindent
Hasonló az összes többi dupla kettős illetve hármas mássalhangzó esete is:
\begin{verbatim}
`ddz `ddzs `ggy `lly `nny `ssz `tty `zzs
`DDZ `DDZS `GGY `LLY `NNY `SSZ `TTY `ZZS
\end{verbatim}

\section{Matematikai képletek}\label{sec-matek}

\begin{latexcode}
\MathBrk{(*\param{matematikai szimbólum}*)}
\MathBrkAll{(*\param{matematikai szimbólum}*)}
\end{latexcode}
A magyar tipográfia a főszövegen belüli belső képletekben a bináris relációk és a bináris operátorok körül (a szorzásjelet és perjelet kivéve) megengedi a sortörést, de sortörés esetén az operátort illetve relációt a következő sor elején meg kell ismételni.
Erre szolgál a \verb|\MathBrk| parancs, amely egyetlen szimbólumot vár paraméterül, ami vagy egyetlen karakter (pl.~\texttt{+}), vagy egy parancs (pl.~\verb|\cup|).
A \verb|\MathBrk| csak akkor elérhető, ha a \textsf{magyar.ldf}-et megfelelő \texttt{mathbrk=} opcióval lett betöltve.
A \verb|\MathBrk| parancs csak akkor módosítja a szimbólumot, ha a magyar nyelv aktív, a \verb|\MathBrkAll| viszont mindig.

Ha a \textsf{magyar.ldf} \texttt{mathbrk=fix} opcióval van betöltve, akkor a szorzásjel és a perjel mellett nem engedélyezett a sortörés, és a legtöbb bináris reláció és bináris operátor a \verb|\MathBrk| kirakása nélkül is a második sorban ismételve választódik el, feltéve, hogy a magyar nyelv aktív.

\begin{latexcode}
\factorial
\end{latexcode}
Ha a \textsf{magyar.ldf}-et a \texttt{mathfactorial=fix} opcióval lett betöltve, akkor a képletekben a felkiáltójel körüli térköz helyes lesz, tehát a faktoriális  egyszerűen bevihető. 
Például \verb|$x=n!k!$| $\to$ $x=n!k!$.
Az eredményben a szorzótényezők közötti szóköz hiányozna \texttt{mathfactorial=fix} nélkül.
Ha csak az alapértelmezés szerinti \texttt{mathfactorial=define} opció aktív, akkor \texttt{\string!} helyett \verb|\factorial|-t érdemes gépelni, hogy a térközök helyesek legyenek.

\begin{latexcode}
\MathReal{(*\param{képlet tizedestörtekkel}*)}
\end{latexcode}
Ezt a parancsot a \texttt{mathreal} opció definiálja.
A paraméterül kapott matematikai képletet úgy jeleníti meg, hogy minden pontot tizedesvesszőre cserél, ha a magyar nyelv aktív.
A paraméterben a pontokat tilos makróba rejteni.
A \verb|\MathReal| képleten belül és azon kívül is működik.
Például

\begin{latexexample}
\MathReal{1.25+2+3.75=7}\\
$\MathReal{1.25+2+3.75=7}$
\end{latexexample}

\section{Automatikus határozott névelő}\label{sec-nevelo}

\begin{latexcode}
\az{(*\param{szó}*)} % a vers 
\Az{(*\param{szó}*)} % Az ős kaján
\az+(*\param{parancs}*){(*\param{argumentum}*)} % az 5. táblázat
\Az+(*\param{parancs}*){(*\param{argumentum}*)} % Az 5. táblázat
\aref{(*\param{címke}*)} % a 2.11. alfejezet
\Aref{(*\param{címke}*)} % A 2.11. alfejezet
\apageref{(*\param{címke}*)} % a 15. oldalon
\Apageref{(*\param{címke}*)} % A 15. oldalon
\acite{(*\param{irodalom-címkék}*)} % az [1, 2] művek rövidek
\Acite{(*\param{irodalom-címkék}*)} % A [3, 4] művek szépek
\end{latexcode}
A \textsf{magyar.ldf} számos parancsot tartalmaz, melyek automatikusan ki tudják tenni a határozott a/az névelőt a szavak elé.
Ezek csak akkor érhetők el, ha a \textsf{magyar.ldf} a megfelelő \texttt{az=} opcióval lett betöltve.
Ezen parancsok csillagos változata csak a névelőt, a csillag nélküli változat pedig a névelővel ellátott szót adja.

Az \verb|\az| és \verb|\Az| parancsok \param{szó} paramétere tartalmazhatja a \verb|\ref|, \verb|\pageref| és \verb|\cite| parancsokat, továbbá a \texttt{+} kiválthat egy kapcsos zárójelet.
Így például a következők ugyanazt szúrják be: \verb|\Az{\pageref{nagy}}|, \verb|\Apageref{nagy}|, \verb|\Az+\pageref{nagy}|.

A \param{szó} elején levő írásjeleket nem veszi figyelembe, például \verb|\az{$*^+$52}| és \verb|\az{52}| ugyanazt eredményezi.
Megzavarhatja a működést, ha a \param{szó}-ban a makrók kifejtése után kapcsos zárójelek vagy kifejthető token-ek maradnak.

A \textsf{magyar.ldf} az arab és római számok és a betűk elé is a megfelelő névelőt rakja.
Például

\begin{latexexample}
\az{xilofon}\\
\az{x betű}\\
\az{\romannumeral10}. osztály\\
\az{\csname @Roman\endcsname{10}}.
osztály\\
\end{latexexample}

A római számokat akkor is helyesen kezeli, ha azok pl.~egy \verb|\aref| által hivatkozott fejezetszám elején találhatók.
Ehhez segítségként az \texttt{.aux} fájlba minden \verb|\newlabel| parancs után a \textsf{magyar.ldf} elhelyez egy \verb|\hunnewlabel| parancsot, amely már mindenképpen arab számokat tartalmaz.
(Ez csak akkor működik, ha a magyar.ldf-et a \texttt{hunnewlabel=yes} opcióval töltöttük be.)

Az előbbi parancsok helyett az 
\begin{latexcode}
\azr \Azr \azp \Azp \azc \Azc 
\end{latexcode}
rövidítések is használhatók, ha a \textsf{magyar.ldf}-et a \texttt{shortrefs=yes} opcióval lett betöltve.

Egyenletekre a következő módokon lehet hivatkozni az \textsf{amsmath} csomag használata esetén:

\begin{latexcode}
\aref({(*\param{címke}*)})
\Aref({(*\param{címke}*)})
\az{\eqref{(*\param{címke}*)}}
\Az{\eqref{(*\param{címke}*)}}
\end{latexcode}
Ezekből inkább az utolsó kettőt ajánlott használni, mert az az egyenletszámot mindig álló betűvel jeleníti meg, ahogy a névelő nélküli esetben is.

A \textsf{magyar.ldf} automatikus határozott névelő kezelésének van néhány korlátja, melyekről a \textsf{huaz} csomag leírásában esik szó (\url{https://ctan.org/pkg/huaz}).
Az automatikus határozott névelő kezeléséhez ez a csomag is használható.

\section{Számok toldalékolása}\label{sec-toldalek}

\begin{latexcode}
\told(*\param{szám}*)+(*\param{toldalék}*){}
\told(*\param{szám}*)+(*\param{toldalék\textsubscript1}*)+(*\param{toldalék\textsubscript2}*){}
\atold(*\param{szám}*)+(*\param{toldalék}*){}
\Atold(*\param{szám}*)+(*\param{toldalék}*){}
\end{latexcode}
Különösen a képletekre való hivatkozásnál lehet hasznos, hogy a \textsf{magyar.ldf} a számneveket ragozni tudja (kivéve, ha ez a \texttt{told=} opcióban le van tiltva).
Így például a ,,(3)-at hozzáadva (4)-hez, majd azt kivonva (6)-ból a (7)-et kapjuk'' alakú kifejezéseket akkor is írhatunk, ha a pontos számokat a \LaTeX{} menet közben generálja.
Ennek a forrása így kezdődik: \verb|\told\aref({eq3})+at{}...|
Ha a dokumentum átrendezése során az \texttt{eq3} címkéjű képlet sorszáma megváltozik, a \verb|\told| akkor is a helyes toldalékot fogja hozzátenni.

A \param{szám} argumentumot az egyértelműség kedvéért kapcsos zárójelbe szabad tenni.
A \verb|\told| a parancs végén található \verb|{}| jelből tudja, hol van vége a toldaléknak: például a \verb|\told1+be{}| és a \verb|\told1+ben{}| között csak így lehet különbséget tenni.
Bizonyos feltételek mellett (például akkor, ha a parancsot pont vagy vessző követi) a \verb|{}| elhagyható, de érdemes a biztonságra törekedve mindig kitenni.
A \param{szám} állhat matematikai módban, és lehet negatív is (pl.~\verb|\told{$-2$}+edik{}| $\to$ $-2$-edik, olvasd: ,,mínusz kettedik''), de vigyázat, mert két mínusz jel kiejti egymást (pl.~\verb|\told{--2}+edik{}| $\to$ --2-odik, de kevesen olvasnák ,,második''-nak).

A \param{szám} tartalmazhat \verb|\ref|, \verb|\aref|, \verb|\Aref|, \verb|\pageref|, \verb|\apageref|, \verb|\Apageref|, \verb|\cite|, \verb|\acite| és \verb|\Acite| parancsokat, de csak csillag és zárójel nélkül.
Például ebben a dokumentációban \aref{subsec-defaults}.~alszakaszra így is lehet hivatkozni:

\begin{latexexample}
\told\ref{subsec-defaults}+as{} 
alszakasz
\end{latexexample}

Az \verb|\atold| illetve \verb|\Atold| parancsok egyenértékűek egy \verb|\az|-ba illetve \verb|\Az|-ba ágyazott \verb|\told|-dal.

A \verb|\told| mindig a legutolsó számot ragozza, tehát például ha az \texttt{ott} címkéjű alalszakasz sorszáma 1.2.3, akkor \verb|\told\ref{ott}+ben{}| így jelenik meg: ,,1.2.3-ban''.

A \param{szám} technikai okokból nem tartalmazhat \verb|\refstruc| parancsot (lásd \aref{sec-hivatkozas}.~szakaszban), mivel az nem expandálható. A megoldás: \verb|\told\refstruc{ott}+ben{}| helyett \verb|\refstruc{ott+ben}|-t kell írni.

A számok végére illeszthető toldalékok a következők:

\begin{longtable}[l]{l|l}
\verb|\told42+a{}| & 42-e\\
\verb|\told42+as{}| & 42-es\\
\verb|\told42+ad{}| & 42-ed\\
\verb|\told42+adik{}| & 42-edik\\
\verb|\told42+an{}| & 42-en\\
\verb|\told42+at{}| & 42-t\\
\verb|\told42+on{}| & 42-en\\
\verb|\told42+a{}| & 42-e\\
\verb|\told42+nal{}| & 42-nél\\
\verb|\told42+ul{}| & 42-ül\\
\verb|\told42+val{}| & 42-vel\\
\verb|\told42+hoz{}| & 42-höz\\
\verb|\told42+ban{}| & 42-ben\\
\verb|\told42+nak{}| & 42-nek\\
\verb|\told42+ba{}| & 42-be\\
\verb|\told42+ra{}| & 42-re\\
\verb|\told42+tol{}| & 42-től\\
\verb|\told42+rol{}| & 42-ről\\
\verb|\told42+szor{}| & 42-ször
\end{longtable}
\noindent
A folytatható toldalékok (\param{toldalék\textsubscript1}) a következők:

\begin{longtable}[l]{l|l}
\verb|\told42+as+at{}| & 42-eset\\
\verb|\told42+ad+szor{}| & 42-edszer\\
\verb|\told42+adik+nal{}| & 42-ediknél\\
\verb|\told42+a+an{}| & 42-én
\end{longtable}
\noindent
Minden \param{toldalék\textsubscript1} minden \param{toldalék\textsubscript2}-vel kombinálható, de nincs értelme az összes kombinációnak.

A számok automatikus toldalékolásához használható még a \textsf{huaz} csomag (\url{https://ctan.org/pkg/huaz}) is.

\section{Hivatkozás a dokumentum strukturális egységeire}\label{sec-hivatkozas}

\begin{latexcode}
\refstruc{(*\param{szám}*)(*\normalfont[*)+(*\param{toldalékok}*)(*\normalfont]*)}
\end{latexcode}
Ez a parancs abban tér el a \verb|\ref| parancstól, hogy nem csak a strukturális egység számszerű megjelölését, hanem a nevét is megjeleníti.
Csak akkor érhető el, ha a \textsf{magyar.ldf} betöltésekor a \texttt{refstruc=} opcióban nem lett letiltva.
Ha a magyar nyelv aktív, akkor ,,1.2.~alszakasz'' formában jeleníti meg a
címet, egyébként ,,subsection 1.2'' formában; az angol elnevezés az egységet nyitó \LaTeX{} parancsszóval egyezik meg.

Ha a magyar nyelv aktív, akkor a \verb|\refstruc| elfogad toldalékokat is,  \aref{sec-toldalek}.~szakaszban leírt módon:

\begin{latexexample}
\refstruc{sec-toldalek+ben}
\end{latexexample}

A -val/-vel ragot nem tudja helyesen hozzákapcsolni a nevekhez.
Csak számozott egységre lehet hivatkozni, tehát például \verb|\section*|-ra nem.
Ha \verb|\az+| van írva a \verb|\refstruc| elé, akkor a megfelelő a/az határozott névelő kerül kiírásra.
Nem használható a \verb|\refstruc| ábrákra, táblázatokra, egyenletekre vagy tételszerű környezetekre való hivatkozásra; ezek ugyanis a dokumentumnak nem strukturális egységei, így például ha egy szakaszban egy ábra sorszáma ,,3.4'', akkor ,,3.4.~szakasz'' generálódna.
Ennek a fogyatékosságnak az az oka, hogy a hivatkozott egység típusa nem áll rendelkezésre az \texttt{.aux} fájlból a \verb|\refstruc| hívásakor.

\section{Dátumok}\label{sec-datum}

A \LaTeX\ \verb|\today| parancsát a dokumentumosztályok többsége átdefiniálja amerikai dátumformátumra, amely a dokumentum fordításának dátumát fogja mutatni.
A \textsf{babel} minden nyelvváltáskor felülírja ezt a parancsot az adott nyelvnek megfelelően.
A dátumgeneráló parancsok használatára néhány példa (melyben a fordítás dátuma \today):

\begin{latexexample}
\SafeToday\\
\today\\
\emitdate{b}{\today}\\
\emitdate[e]{g}{\today}\\
\ondatemagyar\\
\ontoday\\
\emitdate[a+an]{g}{\today}
\end{latexexample}
Ezek a \textsf{magyar.ldf} \texttt{hutoday=} illetve \texttt{emitdate=} opcióival válnak elérhetővé.

Az \verb|\ontoday| parancs csak magyar nyelven elérhető, és az -án/-én raggal látja el a magyar \verb|\today| dátumát.
Az \verb|\ontoday| aktív nyelvtől független, de magyar kimenetet produkáló változata az \verb|\ondatemagyar|.

\begin{latexcode}
\emitdate[(*\param{toldalék}*)]{(*\param{formátum}*)}{(*\param{dátum}*)}
\end{latexcode}
Felismeri a \param{dátum}-ot, és megjeleníti az előírt \param{formátum}-ban, esetleg \param{toldalék}-kal látva el.
A \param{toldalék} a \verb|\told| parancs \texttt{+}-a után megengedett tetszőleges toldalék lehet, tipikus értékei: \texttt{e} és \texttt{a+an}.
Az \texttt{edik+e} használata nem ajánlott, helyette \texttt{e} javasolt.

A \param{dátum} lehet \verb|\today| parancs, vagy egy fix dátum a következő alakokban:
\begin{itemize}[noitemsep,label=]
\item \verb|É-H-N     | (pl.~\texttt{2000-04-01} vagy \texttt{2000-4-1})
\item \verb|É/H/N     | (pl.~\texttt{2000/04/01} vagy \texttt{2000/4/1})
\item \verb|É.H.N     | (pl.~\texttt{2000.04.01} vagy \texttt{2000.4.1})
\item \verb|É. hónap N| (pl.~\texttt{2000. április 1} vagy \texttt{2000. April 1})
\item \verb|hónap N, É| (pl.~\texttt{április 1, 2000} vagy \texttt{April 1, 2000})
\end{itemize}

A 0\dots49 közötti évszámokhoz 2000-et, az 50\dots99 közöttiekhez pedig 1900-at ad hozzá.
A dátum végén lehet egy extra pont, amit figyelmen kívül hagy.
Szóközök számítanak.
Ha a megadott dátum nem felel meg egyik formátumnak sem, az \verb|\emitdate| nem minden esetben ad értelmes hibaüzenetet, és esetleg a dokumentum fordítása végleg elakadhat.

Az \verb|\emitdate| \param{formátum} argumentuma a kimeneti formátumot adja meg.
A magyar helyesírásban minden (\texttt{a}--\texttt{h}) formátum helyes, de \texttt{h} esetén a \textsf{magyar.ldf} nem oldja meg a hónapnevek automatikus ragozását.
A lehetséges formátumok:

\begin{latexexample}
\emitdate{a}{1848-3-15}\\
\emitdate{b}{1848-3-15}\\
\emitdate{c}{1848-3-15}\\
\emitdate{d}{1848-3-15}\\
\emitdate{e}{1848-3-15}\\
\emitdate{f}{1848-3-15} közepe\\
\emitdate[a+an]{g}{1848-3-15}\\
\emitdate{h}{1848-3-15}ában
\end{latexexample}

\section{Lábjegyzetek}\label{sec-labjegyzet}

\subsection{Beállítási lehetőségek}

\begin{latexcode}
\footnotestyle{(*\param{előíráslista}*)}
\end{latexcode}
A \textsf{magyar.ldf} betöltésekor a \texttt{footnote=} opció értékétől függően állítja be a lábjegyzetek stílusát.
\texttt{footnote=unchanged} esetén a lábjegyzetkezelő rész (benne a \verb|\footnotestyle| paranccsal) be se töltődik.
Az alapértelmezés \texttt{footnote=yes}, ekkor a parancsok betöltődnek, de a stílus nem módosul (tehát kezdetben csak a kernel, az osztály és a betöltött csomagok alakítják).
\texttt{footnote=huplain} esetén a \verb|\footnotestyle{huplain}|, illetve \texttt{footnote=starplain} esetén a \verb|\footnotestyle{starplain}| kerül automatikusan kiadásra.

A \verb|\footnotestyle| változtat a lábjegyzetek megjelenítési stílusán.
Hatása csak csoporton belül érvényesül. 
(A lap alján megjelenő részek, például \verb|\footnoterule|, általában az oldal vagy hasáb kiadásakor kerülnek szedésre, ami később történik, mint az oldal főszövegének szedése.
Így például egy oldalon, csoporton belül kiadott \verb|\footnotestyle| visszahathat az előző oldal fejlécét és láblécét alkotó részekre.)
Az \param{előíráslista} vesszővel elválasztott előírásokból áll, melyek sorrendje lényeges.

\medskip
Az előírások a következők lehetnek:

\begin{description}

\begin{footstyle}{reset}

\item[none]
A dokumentumban folyamatosan számozza a lábjegyzeteket.
Csak akkor kezdi újra a számozást, ha egyéb csomagok vagy az osztály ezt állítja be.
Például \textsf{book} osztály esetén a lábjegyzetszámozás fejezetenként újrakezdődik.

\item[section]
Szakaszonként és fejezetenként újrakezdi a lábjegyzetek számozását.

\item[chapter]
Fejezetenként újrakezdi a lábjegyzetek számozását.

\item[page]
1-től azonnal, és a továbbiakban minden oldalon újrakezdi a lábjegyzetek számozását.
Az oldalhatárok követéséhez az \textsf{.aux} fájlt használja, és a konzolon figyelmeztet, ha a változások miatt a \LaTeX-et újra kell futtatni.
Hatóköre végeztével az általa befejezett számozás folytatódik.

\item[page-resume]
1-től azonnal, és a továbbiakban minden oldalon újrakezdi a lábjegyzetek számozását.
Hatóköre végeztével az őt megelőző számozás onnan folytatódik, ahol abbamaradt.

\item[page-cont]
Minden oldalon újrakezdi a lábjegyzetek számozását.
Egy oldalon belül többször ki- és bekapcsolható, emiatt nem kezdi 1-től újra.
A \texttt{page-cont} és a rajta kívüli lábjegyzetek számozása egymástól függetlenül folyamatos.
Az \verb|\editorfootnote| is ezt használja.

\end{footstyle}

\item[resume]
A csoport végén visszaáll az épp aktuális lábjegyzetszám. 
Kiegészítő funkció, önmagában nem használatos.
Például a \verb|\footnotestyle{resume,reset=page}| és a \verb|\footnotestyle{reset=page-resume}| ugyanazt eredményezi.

\begin{footstyle}{indent}

\item[article-nosp]
Az \textsf{article} osztállyal egyező lábjegyzet-megjelenést állít be: a lábjegyzet első sora beljebb kezdődik a többinél, és a sor eleji lábjegyzetszám után nincs szóköz.

\item[article-sp]
Az \textsf{article} osztályhoz hasonló, a magyar tipográfiának megfelelő lábjegyzet`-megjelenést állít be:
a lábjegyzet első sora beljebb kezdődik a többinél, és a sor eleji lábjegyzetszám után egy rövid szóköz van.

\item[hulist]
A magyar tipográfiának megfelelő listaszerű lábjegyzet-megjelenést állít be:
a lábjegyzet első bekezdésének minden sora a bal margótól az első szintű listával egyező mértékben kezdődik beljebb, a lábjegyzetszám ettől balra, egy rövid szóközzel elválasztva lóg ki.

\end{footstyle}

\begin{footstyle}{rule}

\item[none]
Nem helyez el vízszintes vonalat a főszöveg és a lábjegyzetek között.

\item[one-line]
Nem helyez el vízszintes vonalat a főszöveg és a lábjegyzetek között, és e két téglalapot egymástól egy sor $\pm10$\,\% távolságra helyezi el.
A sorméret a \verb|\normalsize| beállítás \verb|\baselineskip| értéke.
Magyar tipográfiában ez a fix távolság ajánlott.
A $\pm10$\,\% esélyt ad a \LaTeX-nek, hogy széthúzza az oldalt.

\item[fourth]
Egy vízszintes vonalat helyez el a főszöveg és a lábjegyzetek között, melynek hossza a szedéstükör egynegyede (tipográfiai szempontból ez ajánlott).
A két téglalap távolságát nem befolyásolja.

\item[choose]
Három különbözőféle vonal közüli választást állít be a főszöveg és a lábjegyzetek között.
A \verb|\pagefootnoterule| makrót fogja hívni ha a legfelső lábjegyzet az aktuális (és nem egy korábbi) oldalon kezdődött.
Alapértéke az eredeti \verb|\footnoterule| vonalhúzómakró.
Ha a legfelső lábjegyzet egy korábbi oldalon kezdődött, akkor az \verb|\mpfootnoterule| makrót fogja hívni, melynek alapértéke a \verb|\footnotestyle{rule=fourth}|-szal egyező.
A \texttt{minipage} környezetben viszont egységesen az \verb|\mpfootnoterule| makrót hívja meg, melynek alapértéke az eredeti \verb|\footnoterule| vonalhúzómakró.

Az előbbi három makrót a szerző a preambulumban átdefiniálhatja.
Arra figyelni kell, hogy a makrókat a \LaTeX{} -- ritka szerencsétlen esetben -- egy oldalon többször is meghívhatja, tehát úgy kell megírni őket, hogy többszöri meghívásuk is ugyanúgy nézzen ki, mint az egyszeri.
Ha a vonalak alapvastagsága 0,4\,pt, akkor például az alábbi definíciók helyesek:\par
\begin{footnotesize}
\begin{verbatim}
\def\pagefootnoterule{\kern-3pt \hrule width 2in \kern2.6pt}
\def\splitfootnoterule{\kern-3pt \hrule width .25\textwidth \kern2.6pt}
\def\mpfootnoterule{}
\end{verbatim}
\end{footnotesize}
Az előbbi három makró elnevezése és használata kompatibilis a \texttt{footmisc} csomag \texttt{splitrule} opcióval történő betöltésével. 

A lábjegyzet és a lábléc közé néha beékelődhet egy úszó objektum.
Ez ellen a lefelé úsztatás tiltásával lehet védekezni, pl.~így: \verb|\begin{table}[th!]|.
Az opciók között ne legyen \texttt{b}, mert az írja elő az oldal alját.

\end{footstyle}

\begin{footstyle}{marksize}

\item[max-normal]
A főszövegben megjelenő lábjegyzetjel méretét korlátozza:
ha az aktuális betűméret nagyobb \verb|\normalsize|-nál, akkor \verb|\footnotesize|-ra korlátozza a jel betűméretét, egyébként nem változtat rajta.

\end{footstyle}

\begin{footstyle}{mark}

\item[arabic]
A lábjegyzetszámok arab megjelenítését állítja be.

\item[stars]
A lábjegyzetszámokat egyre több csillaggal jeleníti meg.
Három csillag fölött már kényelmetlenül hosszú ez a számozás, helyette a \texttt{stars-max} ajánlott.

\item[stars-max]
Az első hét lábjegyzetet így számozza: \textsuperscript{* *** *** † †† ††† *7}. 
Ezután az arab számozás folytatódik, csillaggal kezdve.

\end{footstyle}

\item[mpmark=\param{érték}]
A lehetséges értékek ugyanazok, mint a \texttt{mark=} esetén, és ugyanazt is csinálják, de ezek a a \texttt{minipage} környezeten belüli lábjegyzetjel (\verb|\mpfootnotemark|) megjelenését módosítják.
Továbbá módosítják a \verb|\makeFootnotable| parancs (lásd később) által előkészített környezeteket is.

\item[plain]
Egyenértékű a
\begin{verbatim}
\footnotestyle{reset=none,mark=arabic}
\end{verbatim}
beállítással, ami a \LaTeX{} alapbeállítása.

\item[huplain]
Egyenértékű a
\begin{verbatim}
\footnotestyle{marksize=max-normal,reset=none,
               mark=arabic,rule=none,indent=hulist}
\end{verbatim}
beállítással, ami a \LaTeX{} alapbeállításához hasonló, de a magyar tipográfiát követő arab számos stílus.

\item[starplain]
Egyenértékű a
\begin{verbatim}
\footnotestyle{marksize=max-normal,reset=page,
               mark=stars-max,rule=none,indent=hulist}
\end{verbatim}
beállítással, ami a magyar tipográfiát követő csillagos stílus.

\item[editor]
Egyenértékű a
\begin{verbatim}
\footnotestyle{marksize=max-normal,reset=none,
               mark=arabic,rule=none,indent=hulist}
\end{verbatim}
beállítással, ami a szerkesztő megjegyzéseihez használható.
A szerző lábjegyzeteit ekkor \texttt{huplain} stílusban ajánlott elhelyezni.
Lásd még az \verb|\editorfootnote| parancsot.

\end{description}

\begin{latexcode}
\makeFootnotable{(*\param{környezetnév}*)}
\end{latexcode}
A \param{környezetnév} egy úsztatott környezet neve (pl.~\texttt{figure}).
Ezt az úsztatott környezetet felkészíti a \verb|\footnote| parancs befogadására.
Az úsztatott környezetekben eredetileg a \verb|\footnote| parancs hatására csak a lábjegyzetjel jelenik meg, maga a lábjegyzet hiányozni fog a dokumentumból.
Például a \verb|\makeFootnotable{figure}| kiadása után viszont minden \texttt{figure} környezetben meg fognak jelenni a lábjegyzetek, mégpedig a cím (\verb|\caption|) alatt.
A működés a magyar tipográfiai gyakorlatnak megfelelő, ugyanis az ábra vagy táblázat lábjegyzetét nem a főoldalon, hanem magában az ábrában kell elhelyezni.
A \verb|\makeFootnotable{}| csak az aktuális úsztatott környezet aktív példányára van hatással.

Alternatív megoldás, ha a szóban forgó környezet tartalmát egy \texttt{minipage} környezetben helyezzük el, ekkor azonban számolni kell a \texttt{minipage} egyéb velejáróival is.

Ezt a parancsot a \textsf{footmisc} csomag \texttt{para} opciójával együtt használva, a lábjegyzetek egymás alá fognak kerülni.
Az \textsf{fnpara} csomag esetén viszont e belső lábjegyzetek is egymás mellett jelennek majd meg.

A belső lábjegyzetek megjelenítési stílusát a \verb|\thempfootnote| szabályozza, ami a \LaTeX{} kernel alapértelmezésében dőlt kisbetűket állít be.
Átdefiniálható például félkövér római számokra  \verb|\def\thempfootnote{\bfseries\roman{mpfootnote}}| módon, csillagokra pedig a
\verb|\footnotestyle{mpmark=stars}| paranccsal.

\begin{latexcode}
\editorfootnote{(*\param{lábjegyzet szövege}*)}
\end{latexcode}
A rendes lábjegyzetek közé egy \verb|\footnotestyle{starplain}| stílusban, csillagokkal, oldalanként újrakezdődő számozású lábjegyzetet helyez el.
Hatása megegyezik egy csoporton belül \verb|\footnotestyle{editor}| után kiadott \verb|\footnote|-éval.
Jól használható a szerkesztő megjegyzéseinek feltüntetésére, miközben a szerző jegyzetei sima \verb|\footnote| paranccsal szerepelnek.

\begin{latexcode}
\setlength{\skip\footins}{(*\param{rugalmas távolság}*)}
\setlength{\dimen\footins}{(*\param{maximális magasság}*)}
\end{latexcode}
A főszöveg és a lábjegyzettömb távolságát, illetve a lábjegyzettömb maximális magasságát állítja be.
Előbbi alapértéke \verb|\bigskipamount|, utóbbié pedig \texttt{8in}, de ezt az osztály és a betöltött csomagok felülbírálhatják.

\begin{latexcode}
\headingfootnote(*\textrm{[}[\dots]\textrm{]}*){(*\param{lábjegyzet szövege}*)}
\end{latexcode}
Fejezet- illetve szakaszcímekben nem lehet használni a \verb|\footnote| parancsot.
Ha mégis szükséges, akkor \verb|\footnote| helyett a \verb|\headingfootnote| parancsot kell használni.

\subsection{A lábjegyzetekre vonatkozó magyar tipográfiai szabályok}

A lábjegyzetek magyar tipográfiai szabályainak a \textsf{magyar.ldf} a következő listában leírt módokon próbál megfelelni:

\begin{enumerate}
\item
Csillagok és arab számozás egyaránt használható a lábjegyzet jelölésére.
A kívánt számozás a \texttt{mark=} előírással kiválasztható.
\item
A csillagos számozás minden oldalon újrakezdődik, az arab számozás pedig fejezetenként.
Lásd a \texttt{reset=} előírást.
\item
A lábjegyzetjel szóköz nélkül követi a szót, amire vonatkozik.
Ennek érdekében a szerzőnek kell figyelni arra, hogy a \verb|\footnote| előtt ne legyen szóköz vagy sortörés.
\item
A lábjegyzet külön mondat, tehát nagybetűvel kell kezdeni, és általában ponttal kell befejezni.
Erre a szerzőnek kell figyelni.
\item
Ha egy műben a szerkesztő és a szerző is helyez el lábjegyzeteket, a szerző lábjegyzeteit arab számokkal, a szerkesztőjét pedig csillagokkal érdemes jelölni. Lásd az \verb|\editorfootnote| parancsot.
\item
A lábjegyzetek szövege és a főszöveg között elég egy üres sort hagyni, vonalra nincs
szükség.
Lásd a \texttt{rule=none} vagy \texttt{rule=one-line} előírásokat.
\item
A főszöveg és egy hosszú lábjegyzet előző oldalon félbemaradt részének folytatásával induló lábjegyzetblokk között vonalat kell elhelyezni (melynek hossza a tükörszélesség harmada vagy negyede legyen).
A legtöbbször jól működő megoldást valósít meg a \texttt{rule=choose} előírás.
\item
Rövid lábjegyzetek egy sorba tördelhetők.
Ehhez a \textsf{footmisc} csomag ajánlott \texttt{para} opcióval, vagy az \textsf{fnpara} csomag. 
A \textsf{footmisc} a \texttt{minipage} környezeten belüli lábjegyzeteket nem teszi egy sorba, az \textsf{fnpara} viszont igen.
\item
Táblázat illetve ábra lábjegyzetei közvetlenül alatta, és ne az oldal alján jelenjenek meg.
Lásd a \verb|\makeFootnotable| parancsot.
\item
Ha a normálnál nagyobb betűméretű szövegben van lábjegyzet, a lábjegyzetjel ne legyen nagyobb a normál betűmérethez tartozónál.
Lásd a \texttt{marksize=max-normal} előírást.
\item
A lábjegyzet kezdődhet beütéssel (azaz első sora beljebb kezdődik) vagy listaszerűen (a lábjegyzetjeltől eltekintve az első és a második sor egymás fölött kezdődik).
Lásd az \texttt{indent=hulist} előírást.
\item
Az oldal alján található lábjegyzetjel után egy rövid szóköz van.
Ehhez használhatóak az \texttt{indent=article-sp} vagy \texttt{indent=hulist} előírások.
\item
Ha túl sok csillagos megjegyzés kerülne egy oldalra, akkor más jelre kell áttérni.
Lásd a \texttt{mark=stars-max} előírást.
\end{enumerate}

\end{document}